\documentclass[11pt,a4paper]{article}
\usepackage[margin=1.8cm]{geometry}
\usepackage{amsmath,amssymb,array,tabularx,fontspec,cancel,tikz,graphicx}
\setmainfont{FreeSerif}
\setlength{\parindent}{0pt}
\setlength{\parskip}{0pt}
\newcounter{sq}
% \sq{text}: auto-numbered (Roman numerals like the original paper)
\newcommand{\sq}[1]{\stepcounter{sq}\par\medskip\noindent\hangindent=7em\hangafter=1
  \makebox[7em][l]{\textbf{\Roman{sq}.}}#1\par}

% \fd{left labels}{right labels}{edges}: arrow (mapping) diagram
\newcommand{\fd}[3]{\par\smallskip\noindent\hspace*{7em}\begin{tikzpicture}[>=stealth,scale=0.75]
 \draw (0,0) ellipse (0.9 and 1.5); \draw (3.5,0) ellipse (0.9 and 1.5);
 \foreach \l [count=\i] in {#1} {\node[font=\small] (L\i) at (0,1.25-0.7*\i+0.7) {\l};}
 \foreach \r [count=\i] in {#2} {\node[font=\small] (R\i) at (3.5,1.25-0.7*\i+0.7) {\r};}
 #3
 \end{tikzpicture}\par}
% \tfig{tikz code}: geometry figure
\newcommand{\tfig}[1]{\par\smallskip\noindent\hspace*{7em}\begin{tikzpicture}[scale=0.8,font=\small]#1\end{tikzpicture}\par}
% \pfig{file}: figure copied (cropped) from the original paper
\newcommand{\pfig}[2]{\par\smallskip\noindent\hspace*{7em}\includegraphics[width=#2]{#1}\par}
% score-distribution table used in several questions
\newcommand{\dist}{\par\smallskip\noindent\hspace*{7em}\begin{tabular}{|l|c|c|c|c|c|c|c|}\hline
\textbf{Scores} & 24-28 & 29-33 & 34-38 & 39-43 & 44-48 & 49-53 & \textbf{Total}\\ \hline
\textbf{No. of students} & 3 & 6 & 12 & 23 & 15 & 6 & 65\\ \hline\end{tabular}\par\smallskip\hspace*{7em}}
% biology marks tables (6 columns / 5 columns as in the original)
\newcommand{\bioA}{\par\smallskip\noindent\hspace*{7em}\begin{tabular}{|l|c|c|c|c|c|c|}\hline
\textbf{Marks} & 20-24 & 25-29 & 30-34 & 35-39 & 40-44 & 45-49\\ \hline
\textbf{No. of students} & 05 & 08 & 12 & 15 & 03 & 02\\ \hline\end{tabular}\par\smallskip\hspace*{7em}}
\newcommand{\bioB}{\par\smallskip\noindent\hspace*{7em}\begin{tabular}{|l|c|c|c|c|c|}\hline
\textbf{Marks} & 20-24 & 25-29 & 30-34 & 35-39 & 40-44\\ \hline
\textbf{No. of students} & 05 & 08 & 12 & 15 & 03\\ \hline\end{tabular}\par\smallskip\hspace*{7em}}

\begin{document}
\begin{center}
  {\Huge\bfseries The Hope Science Academy}\\[3pt]
  {\small Excellence is our standard not goal}\\[2pt]
  {\small 247-E-1, (Opposite LDA School) Johar Town, Lahore. PH: 0300-8194789, 0322-4157001}
\end{center}
\medskip\noindent
\begin{tabularx}{\linewidth}{|X|X|X|X|}
\hline
\textbf{Student Name:} & \textbf{Roll Num:} & \textbf{Class:} 9th & \textbf{Paper Code:} \\ \hline
\textbf{Subject Name:} Mathematics. & \textbf{Time Allowed:} & \textbf{Total Marks:} 1374 & \textbf{Exam Date:} 09-Oct-2026 \\ \hline
\multicolumn{4}{|l|}{\textbf{Exam Syllabus:} FULL BOOK} \\ \hline
\end{tabularx}
\bigskip
\noindent\textbf{Q2. Write short answers of the following questions.}\hfill\textbf{1$\times$831=831}

\sq{Identify the following as a rational or irrational number. 2.353535}
\sq{Identify the following as a rational or irrational number. $0.\overline{6}$}
\sq{Identify the following as a rational or irrational number. 2.236067 \ldots}
\sq{Identify the following as a rational or irrational number. $\sqrt{7}$}
\sq{Identify the following as a rational or irrational number. e}
\sq{Identify the following as a rational or irrational number. $\pi$}
\sq{Identify the following as a rational or irrational number. $5+\sqrt{11}$}
\sq{Identify the following as a rational or irrational number. $\sqrt{3}+\sqrt{13}$}
\sq{Identify the following as a rational or irrational number. $\dfrac{15}{4}$}
\sq{Identify the following as a rational or irrational number. $(2-\sqrt{2})(2+\sqrt{2})$}
\sq{Represent the following number on number line: $\sqrt{2}$}
\sq{Represent the following number on number line: $\sqrt{3}$}
\sq{Represent the following number on number line: $4\dfrac{1}{3}$}
\sq{Represent the following number on number line: $-2\dfrac{1}{7}$}
\sq{Represent the following number on number line: $\dfrac{5}{8}$}
\sq{Represent the following number on number line: $2\dfrac{3}{4}$}
\sq{Express the following as a rational number $\dfrac{p}{q}$ where p and q are integers and $q\neq 0$: $0.\overline{4}$}
\sq{Express the following as a rational number $\dfrac{p}{q}$ where p and q are integers and $q\neq 0$: $0.\overline{37}$}
\sq{Express the following as a rational number $\dfrac{p}{q}$ where p and q are integers and $q\neq 0$: $0.\overline{21}$}
\sq{Name the property used in the following: $(a+4)+b=a+(4+b)$}
\sq{Name the property used in the following: $\sqrt{2}+\sqrt{3}=\sqrt{3}+\sqrt{2}$}
\sq{Name the property used in the following: x -x = 0}
\sq{Name the property used in the following: $a(b+c)=ab+ac$}
\sq{Name the property used in the following: 16 + 0 = 16}
\sq{Name the property used in the following: 100 x 1 = 100}
\sq{Name the property used in the following: 4 x (5 x 8) = (4 x 5) x 8}
\sq{Name the property used in the following: ab=ba}
\sq{Name the property used in the following: $-3<-1\Rightarrow 0<2$}
\sq{Name the property used in the following: If $a<b$ then $\dfrac{1}{a}>\dfrac{1}{b}$}
\sq{Name the property used in the following: If $a<b$ then $a+c<b+c$}
\sq{Name the property used in the following: If $ac<bc$ and $c>0$ then $a<b$}
\sq{Name the property used in the following: if $ac<bc$ and $c<0$ then $a>b$}
\sq{Name the property used in the following: Either $a>b$ or $a=b$ or $a<b$}
\sq{Insert two rational numbers between: $\dfrac{1}{3}$ and $\dfrac{1}{4}$}
\sq{Insert two rational numbers between: 3 and 4}
\sq{Insert two rational numbers between: $\dfrac{3}{5}$ and $\dfrac{4}{5}$}
\sq{Define terminating decimal numbers.}
\sq{Define Non-terminating decimal numbers.}
\sq{Define rational numbers.}
\sq{define irrational numbers}
\sq{Identify the following decimal number as rational or irrational number: 0.35}
\sq{Identify the following decimal number as rational or irrational number: 0.444...}
\sq{Identify the following decimal number as rational or irrational number: $3.\overline{5}$}
\sq{Identify the following decimal number as rational or irrational number: 3.36788542...}
\sq{Identify the following decimal number as rational or irrational number: 1.709975947}
\sq{Represent $\sqrt{5}$ on a number line.}
\sq{Express the following recurring decimal as the rational number $\dfrac{p}{q}$, where p and q are integers. $0.\overline{5}$}
\sq{Express the following recurring decimal as the rational number $\dfrac{p}{q}$, where p and q are integers. $0.\overline{93}$}
\sq{Insert two rational numbers between 2 and 3.}
\sq{Identify the property that justifies the statement: If $a>13$ then $a+2>15$}
\sq{Identify the property that justifies the statement: If $3<9$ and $6<12$ then $9<21$}
\sq{Identify the property that justifies the statement: If $7>4$ and $5>3$ then $35>12$}
\sq{Identify the property that justifies the statement: If $-5<-4\Rightarrow 20>16$}
\sq{Rationalize the denominator of following: $\dfrac{13}{4+\sqrt{3}}$}
\sq{Rationalize the denominator of following: $\dfrac{\sqrt{2}+\sqrt{5}}{\sqrt{3}}$}
\sq{Rationalize the denominator of following: $\dfrac{\sqrt{2}-1}{\sqrt{5}}$}
\sq{Rationalize the denominator of following: $\dfrac{6-4\sqrt{2}}{6+4\sqrt{2}}$}
\sq{Rationalize the denominator of following: $\dfrac{\sqrt{3}-\sqrt{2}}{\sqrt{3}+\sqrt{2}}$}
\sq{Rationalize the denominator of following: $\dfrac{4\sqrt{3}}{\sqrt{7}+\sqrt{5}}$}
\sq{Simplify the following: $\left(\dfrac{81}{16}\right)^{-\frac{3}{4}}$}
\sq{Simplify the following: $\left(\dfrac{3}{4}\right)^{-2}\div\left(\dfrac{4}{9}\right)^{3}\times\dfrac{16}{27}$}
\sq{Simplify the following: $(0.027)^{-\frac{1}{3}}$}
\sq{Simplify the following: $\sqrt[7]{\dfrac{x^{14}\times y^{21}\times z^{35}}{y^{14}z^{7}}}$}
\sq{Simplify the following: $(64)^{-\frac{2}{3}}\div(9)^{-\frac{3}{2}}$}
\sq{If $x=3+\sqrt{8}$ then find the value of: $x+\dfrac{1}{x}$}
\sq{If $x=3+\sqrt{8}$ then find the value of: $x-\dfrac{1}{x}$}
\sq{If $x=3+\sqrt{8})$ then find the value of: $x^{2}+\dfrac{1}{x^{2}}$}
\sq{If $x=3+\sqrt{8}$ then find the value of: $\left(x-\dfrac{1}{x}\right)^{2}$}
\sq{Define radicand.}
\sq{Define Surds.}
\sq{Define monomial surd.}
\sq{Define binomial surds.}
\sq{Define canjugate surds.}
\sq{Is 0 a rational number? Explain.}
\sq{State trichotomy property of real numbers.}
\sq{Find two rational numbers between 4 and 5.}
\sq{If $a=\dfrac{3}{2},b=\dfrac{5}{3}$ and $c=\dfrac{7}{5}$, then verify that: $a(b+c)=ab+ac$}
\sq{If $a=\dfrac{3}{2},b=\dfrac{5}{3}$ and $c=\dfrac{7}{5}$, then verify that: $(a+b)c=ac+bc$}
\sq{If $a=\dfrac{4}{3},b=\dfrac{5}{2},c=\dfrac{7}{4}$, then verify the associative property of real numbers w.r.t addition and multiplicatio}
\sq{State trichotomy property of real numbers.}
\sq{Simplify the following: $\sqrt[5]{\dfrac{x^{15}y^{35}}{z^{20}}}$}
\sq{Simplify the following: $\sqrt[3]{(27)^{2x}}$}
\sq{Simplify the following: $\dfrac{6(3)^{n+2}}{3^{n+1}-3^{n}}$}
\sq{Simplify the following: $\sqrt[4]{16x^{4}y^{8}}$}
\sq{Simplify the following: $\sqrt[3]{27x^{6}y^{9}z^{3}}$}
\sq{Simplify the following: $(64)^{-\frac{4}{3}}$}
\sq{Rationalize the denominator of: $\dfrac{3}{\sqrt{5}+\sqrt{2}}$}
\sq{Rationalize the denominator of: $\dfrac{3}{\sqrt{5}-\sqrt{3}}$}
\sq{The sum of three consecutive integers is forty-two, find the three integers.}
\sq{The sum of two real numbers is 8, and their difference is 2. Find the numbers.}
\sq{Normal human body temperature is 98.6 °F. Convert it into Celsius and Kelvin scale.}
\sq{Express the following number in scientific notation. 2000000}
\sq{Express the following number in scientific notation. 48900}
\sq{Express the following number in scientific notation. 0.0042}
\sq{Express the following number in scientific notation. 0.0000009}
\sq{Express the following number in scientific notation. $73\times10^{3}$}
\sq{Express the following number in scientific notation. $0.65\times10^{2}$}
\sq{Express the following number in ordinary notation. $8.04\times10^{2}$}
\sq{Express the following number in ordinary notation. $3\times10^{5}$}
\sq{Express the following number in ordinary notation. $1.5\times10^{-2}$}
\sq{Express the following number in ordinary notation. $1.77\times10^{7}$}
\sq{Express the following number in ordinary notation. $5.5\times10^{-6}$}
\sq{Express the following number in ordinary notation. $4\times10^{-5}$}
\sq{The speed of light is approximately $3\times10^{8}$ metres per second. Express it in standard form.}
\sq{The circumference of the Earth at the equator is about 40075000 metres. Express this number in scientific notation.}
\sq{The diameter of Mars is $6.7779\times10^{3}$ km. Express this number in standard form.}
\sq{The diameter of Earth is $1.2756\times10^{4}$ km. Express this number in standard form.}
\sq{Convert $5.63\times10^{3}$ to ordinary notation.}
\sq{Convert $6.6\times10^{-5}$ to ordinary notation.}
\sq{Convert 29,000,000 to scientific notation.}
\sq{Convert 0.000006 to scientific notation.}
\sq{Define Scientific notation.}
\sq{Where Scientific notation is used?}
\sq{Convert 78,000,000 to scientific notation.}
\sq{Convert 0.0000000315 to scientific notation.}
\sq{Convert $3.47\times10^{6}$ to ordinary notation.}
\sq{Convert $6.23\times10^{-4}$ to ordinary notation.}
\sq{Express the following in logarithmic form: $10^{3}=1000$}
\sq{Express the following in logarithmic form: $2^{8}=256$}
\sq{Express the following in logarithmic form: $3^{-3}=\dfrac{1}{27}$}
\sq{Express the following in logarithmic form: $20^{2}=400$}
\sq{Express the following in logarithmic form: $16^{-\frac{1}{4}}=\dfrac{1}{2}$}
\sq{Express the following in logarithmic form: $11^{2}=121$}
\sq{Express the following in logarithmic form: $p=q^{r}$}
\sq{Express the following in logarithmic form: $(32)^{\frac{-1}{5}}=\dfrac{1}{2}$}
\sq{Express the following in exponential form: $\log_{5}125=3$}
\sq{Express the following in exponential form: $\log_{2}16=4$}
\sq{Express the following in exponential form: $\log_{23}1=0$}
\sq{Express the following in exponential form: $\log_{5}5=1$}
\sq{Express the following in exponential form: $\log_{2}\dfrac{1}{8}=-3$}
\sq{Express the following in exponential form: $\dfrac{1}{2}=\log_{9}3$}
\sq{Express the following in exponential form: $5=\log_{10}100000$}
\sq{Express the following in exponential form: $\log_{4}\dfrac{1}{16}=-2$}
\sq{Find the value of x in the following: $\log_{x}64=3$}
\sq{Find the value of x in the following: $\log_{5}1=x$}
\sq{Find the value of x in the following: $\log_{x}8=1$}
\sq{Find the value of x in the following: $\log_{10}x=-3$}
\sq{Define logarithm?}
\sq{What do you about logarithm of a real number?}
\sq{Convert $\log_{2}8=3$ to exponential form.}
\sq{Convert $\log_{10}100=2$}
\sq{Find the value of x in the following case: $\log_{5}25=x$}
\sq{Find the value of x in the following case: $\log_{2}x=6$}
\sq{Find the value of x in the following case: $\log_{5}25=x$}
\sq{Covert the following in logarithmic form: $3^{4}=81$}
\sq{Convert the following in logarithmic form: $7^{0}=1$}
\sq{Find characteristic of the following number: 5287}
\sq{Find characteristic of the following number: 59.28}
\sq{Find characteristic of the following number: 0.0567}
\sq{Find characteristic of the following number: 234.7}
\sq{Find characteristic of the following number: 0.000049}
\sq{Find characteristic of the following number: 145000}
\sq{Find logarithm of the following number: 43}
\sq{Find logarithm of the following number: 579}
\sq{Find logarithm of the following number: 1.982}
\sq{Find logarithm of the following number: 0.0876}
\sq{Find logarithm of the following number: 0.047}
\sq{Find logarithm of the following number: 0.000354}
\sq{Find the value of x: log x = 0.0065}
\sq{Find the value of x: log x = 1.192}
\sq{Define common logarith.}
\sq{Define natural log.}
\sq{Define characteristics.}
\sq{Define mantissa.}
\sq{Write rules for finding the characteristics.}
\sq{Deinfe anti-logarithm.}
\sq{Write difference between common and natural logaritm.}
\sq{Without using calculator, evaluate the following: $\log_{2}18-\log_{2}9$}
\sq{Without using calculator, evaluate the following: $\log_{2}64+\log_{2}2$}
\sq{Without using calculator, evaluate the following: $2\log2+\log25$}
\sq{Write the following as a single logarithm: $\dfrac{1}{2}\log25+2\log3$}
\sq{Write the following as a single logarithm: $\log9-\log\dfrac{1}{3}$}
\sq{Write the following as a single logarithm: $\log_{5}b^{2}.\log_{a}5^{3}$}
\sq{Write the following as a single logarithm: $2\log_{3}x+\log_{3}y$}
\sq{Write the following as a single logarithm: $4\log_{5}x-\log_{5}y+\log_{5}z$}
\sq{Write the following as a single logarithm: $3\ln a+3\ln b-\ln c$}
\sq{Expand the following using laws of logarithm: $\log\!\left(\dfrac{11}{5}\right)$}
\sq{Expand the following using laws of logarithm: $\log_{5}\sqrt{8a^{6}}$}
\sq{Expand the following using laws of logarithm: $\ln\!\left(\dfrac{a^{2}b}{c}\right)$}
\sq{Expand the following using laws of logarithm: $\log\!\left(\dfrac{xy}{z}\right)^{\frac{1}{9}}$}
\sq{Expand the following using laws of logarithm: $\ln\sqrt[3]{16x^{3}}$}
\sq{Expand the following using laws of logarithm: $\log_{2}\!\left(\dfrac{1-a}{b}\right)^{5}$}
\sq{Find the value of x in the following equation: $\log2+\log x=1$}
\sq{Find the value of x in the following equation: $\log_{2}x+\log_{2}8=5$}
\sq{Find the value of x in the following equation: $\log(5x-10)=2$}
\sq{Write laws of logarithm.}
\sq{What is change of base law?}
\sq{Prove the change of base law.}
\sq{Prove the product law.}
\sq{Prove the quotient law.}
\sq{Prove the power law.}
\sq{Expand the following using laws of logarithms: $\log_{3}(20)$}
\sq{Expand the following using laws of logarithms: $\log_{2}(9)^{5}$}
\sq{Expand the following using laws of logarithms: $\log_{32}27$}
\sq{Expand the following using laws of logarithms: $\log_{2}\!\left(\dfrac{x-y}{z}\right)^{3}$}
\sq{Expand the following using laws of logarithms: $\log_{5}\!\left(\dfrac{xy}{z}\right)^{8}$}
\sq{Write the following as a single logarithm: $2\log_{3}10-\log_{3}4$}
\sq{Write the following as a single logarithm: $6\log_{3}x+2\log_{3}11$}
\sq{Express the following number in scientific notation: 0.000567}
\sq{Express the following number in scientific notation: 734}
\sq{Express the following number in scientific notation: $0.33\times10^{3}$}
\sq{Express the following number in ordinary notation: $2.6\times10^{3}$}
\sq{Express the following number in ordinary notation: $8.794\times10^{-4}$}
\sq{Express the following number in ordinary notation: $6\times10^{-6}$}
\sq{Express the following in logarithmic form: $3^{7}=2187$}
\sq{Express the following in logarithmic form: $a^{b}=c$}
\sq{Express the following in logarithmic form: $(12)^{2}=144$}
\sq{Express the following in exponential form: $\log_{4}8=x$}
\sq{Express the following in exponential form: $\log_{9}729=3$}
\sq{Express the following in exponential form: $\log_{4}1024=5$}
\sq{Find value of x in the following: $\log_{9}x=0.5$}
\sq{Find value of x in the following: $\left(\dfrac{1}{9}\right)^{3x}=27$}
\sq{Write the following as a single logarithm: $7\log x-3\log y^{2}$}
\sq{Write the following as a single logarithm: $3\log4-\log32$}
\sq{Write the following as a single logarithm: $\dfrac{1}{3}(\log_{5}8+\log_{5}27)-\log_{5}3$}
\sq{Expand the following using laws of logarithms: $\log(x\,y\,z^{6})$}
\sq{Expand the following using laws of logarithms: $\log_{3}\sqrt[6]{m^{5}n^{3}}$}
\sq{Expand the following using laws of logarithms: $\log\sqrt{8x^{3}}$}
\sq{Find the values of the following with the help of logarithm table: $\sqrt[3]{68.24}$}
\sq{Write the following set in set builder notation: \{1,4,9,16,25,36,...,484\}}
\sq{Write the following set in set builder notation: \{2,4,8,16,32,64,...,156\}}
\sq{Write the following set in set builder notation: $\{0,\pm1,\pm2,\ldots,\pm1000\}$}
\sq{Write the following set in set builder notation: \{6,12,18,...,120\}}
\sq{Write the following set in set builder notation: \{100,102,104,...,400\}}
\sq{Write the following set in set builder notation: \{1,3,9,27,81,...\}}
\sq{Write the following set in set builder notation: \{1,2,4,5,10,20,25,50,100\}}
\sq{Write the following set in set builder notation: \{5,10,15,...,100\}}
\sq{Write the following set in set builder notation: The set of all integers between -100 and 1000}
\sq{Write the following set in tabular form: $\{x|x \text{ is a multiple of } 3\wedge\ \leq 35\}$}
\sq{Write the following set in tabular form: $\{x|x\in R\wedge 2x+1=0\}$}
\sq{Write the following set in tabular form: $\{x|x\in P\wedge x<12\}$}
\sq{Write the following set in tabular form: $\{x|x \text{ is a divisor of } 128\}$}
\sq{Write the following set in tabular form: $\{x|x=2^{n},n\in N\wedge n<8\}$}
\sq{Write the following set in tabular form: $\{x|x\in N\wedge x+4=0\}$}
\sq{Write the following set in tabular form: $\{x|x\in N\wedge x=x\}$}
\sq{Write the following set in tabular form: $\{x|x\in Z\wedge 3x+1=0\}$}
\sq{Write two proper subset of the following set: \{a,b,c\}}
\sq{Write two proper subset of the following set: \{0,1\}}
\sq{Write two proper subset of the following set: N}
\sq{Write two proper subset of the following set: Z}
\sq{Write two proper subset of the following set: Q}
\sq{Write two proper subset of the following set: R}
\sq{Write two proper subset of the following set: $\{x\in Q\wedge 0<x\leq2\}$}
\sq{Is there any set which has no proper subset? if so, name that set.}
\sq{What is the difference between \{a,b\} and $\{\{a,b\}\}$?}
\sq{What is the number of element of the power set of each of the following set? \{ \}}
\sq{What is the number of element of the power set of each of the following set? \{0,1\}}
\sq{What is the number of element of the power set of each of the following set? \{1,2,3,4,5,6,7\}}
\sq{What is the number of element of the power set of each of the following set? \{0,1,2,3,4,5,6,7\}}
\sq{What is the number of element of the power set of each of the following set? $\{a,\{b,c\}\}$}
\sq{What is the number of element of the power set of each of the following set? $\{\{a,b\},\{b,c\},\{d,e\}\}$}
\sq{Write down the power set of the following set: \{9,11\}}
\sq{Write down the power set of the following set: $\{+,-,\times,\div\}$}
\sq{Write down the power set of the following set: $\{\phi\}$}
\sq{Write down the power set of the following set: $\{a,\{b,c\}\}$}
\sq{Define a set.}
\sq{What is descriptive method?}
\sq{What is tabular method?}
\sq{What is set-builder method?}
\sq{What is sigleton set?}
\sq{Define Null set.}
\sq{How empty set is written?}
\sq{Define equal set.}
\sq{Define equivalent set.}
\sq{Define subset.}
\sq{Define proper subset.}
\sq{Define improper subset.}
\sq{Define universal set.}
\sq{Define power set}
\sq{Consider the universal set $U=\{x: x \text{ is multiple of } 2 \text{ and } 0<x\leq30\}$, A=\{x;x is a multiple of 6\} and B=\{x:x is a multiple of 8\}: List all elements of sets A and B in tabular form.}
\sq{Consider the universal set $U=\{x: x \text{ is multiple of } 2 \text{ and } 0<x\leq30\}$, A=\{x;x is a multiple of 6\} and B=\{x:x is a multiple of 8\}: Find $A\cap B$}
\sq{If U=\{1,2,3,...,20\} and A=\{1,3,5,...,19\}, verify the following: $A\cap U=A$}
\sq{What is union of two sets?}
\sq{What is intersection of two sets?}
\sq{Define disjoint sets.}
\sq{Define overlapping sets.}
\sq{What is complement of a set?}
\sq{What is Venn Diagram?}
\sq{Who used the Venn Diagram first?}
\sq{What is inclusion-exclusion principle?}
\sq{Define cardinality of a set.}
\sq{For A = \{1,2,3,4\}, find the following relation in A. State the domain and range of the following relation: $\{(x,y)\mid y=x\}$}
\sq{For A = \{1,2,3,4\}, find the following relation in A. State the domain and range of the following relation: $\{(x,y)\mid y+x=5\}$}
\sq{For A = \{1,2,3,4\}, find the following relation in A. State the domain and range of the following relation: $\{(x,y)\mid x+y<5\}$}
\sq{For A = \{1,2,3,4\}, find the following relation in A. State the domain and range of the following relation: $\{(x,y)\mid x+y>5\}$}
\sq{Is the following diagram represent the functions and of which type?
\fd{1,2,3}{a,b,c,d}{\draw[->](L1)--(R1); \draw[->](L1)--(R2); \draw[<->](L2)--(R3); \draw[<->](L3)--(R4);}}
\sq{Is the following diagram represent the functions and of which type?
\fd{a,b,c}{1,3,5}{\draw[->](L1)--(R1); \draw[->](L2)--(R2); \draw[->](L3)--(R3);}}
\sq{Is the following diagram represent the functions and of which type?
\fd{1,2,3}{a,b,c}{\draw[->](L1)--(R1); \draw[->](L2)--(R2); \draw[->](L3)--(R3);}}
\sq{Is the following diagram represent the functions and of which type?
\fd{l,m,n}{x,y,z}{\draw[->](L1)--(R1); \draw[->](L2)--(R1); \draw[->](L3)--(R3);}}
\sq{If $g(x)=3x+2$ then find: g(0)}
\sq{If $g(x)=3x+2$ then find: g(-3)}
\sq{If $g(x)=3x+2$ then find: $g\!\left(\dfrac{2}{3}\right)$}
\sq{If $h(x)=x^{2}+1$, then find: h(1)}
\sq{If $h(x)=x^{2}+1$, then find: h(-4)}
\sq{If $h(x)=x^{2}+1$, then find: $h\!\left(-\dfrac{1}{2}\right)$}
\sq{Consider the function defined by $f(x)=5x+1$. If $f(x)=32$, find the x value.}
\sq{Define functions.}
\sq{How many types of fucnction? write their names.}
\sq{What is Into function?}
\sq{What is Onto function?}
\sq{What is One-One function?}
\sq{What is Bijective function?}
\sq{Define linear function.}
\sq{Define quadratic function.}
\sq{What is binary relation?}
\sq{Define domain.}
\sq{Define range.}
\sq{Write the following set in tabular forms: $\{x\mid x=2n,\ n\in N\}$}
\sq{Write the following set in tabular forms: $\{x\mid x=2m+1,m\in N\}$}
\sq{Write the following set in tabular forms: $\{x\mid x=11n,n\in W\wedge n<11\}$}
\sq{Write the following set in tabular forms: $\{x\mid x\in E\wedge 4<x<6\}$}
\sq{Write the following set in tabular forms: $\{x\mid x\in O\wedge 5\leq x<7\}$}
\sq{Write the following set in tabular forms: $\{x\mid x\in Q\wedge x^{2}=2\}$}
\sq{Write the following set in tabular forms: $\{x\mid x\in Q\wedge x=-x\}$}
\sq{Write the following set in tabular forms: $\{x\mid x\in R\wedge x\notin Q'\}$}
\sq{Let $U=\{1,2,3,4,5,6,7,8,9,10\}$, $A=\{2,4,6,8,10\}$, $B=\{1,2,3,4,5\}$ and $C=\{1,3,5,7,9\}$ List the member of the following set: $A'$}
\sq{Let $U=\{1,2,3,4,5,6,7,8,9,10\}$, $A=\{2,4,6,8,10\}$, $B=\{1,2,3,4,5\}$ and $C=\{1,3,5,7,9\}$ List the member of the following set: $B'$}
\sq{Let $U=\{1,2,3,4,5,6,7,8,9,10\}$, $A=\{2,4,6,8,10\}$, $B=\{1,2,3,4,5\}$ and $C=\{1,3,5,7,9\}$ List the member of the following set: $A\cup B$}
\sq{Let $U=\{1,2,3,4,5,6,7,8,9,10\}$, $A=\{2,4,6,8,10\}$, $B=\{1,2,3,4,5\}$ and $C=\{1,3,5,7,9\}$ List the member of the following set: A-B}
\sq{Let $U=\{1,2,3,4,5,6,7,8,9,10\}$, $A=\{2,4,6,8,10\}$, $B=\{1,2,3,4,5\}$ and $C=\{1,3,5,7,9\}$ List the member of the following set: $A\cap C$}
\sq{Let $U=\{1,2,3,4,5,6,7,8,9,10\}$, $A=\{2,4,6,8,10\}$, $B=\{1,2,3,4,5\}$ and $C=\{1,3,5,7,9\}$ List the member of the following set: $U'$}
\sq{Using the Venn diagrams, if necessary, find the single sets equal to the following: $A'$}
\sq{Using the Venn diagrams, if necessary, find the single sets equal to the following: $A\cap U$}
\sq{Using the Venn diagrams, if necessary, find the single sets equal to the following: $A\cup U$}
\sq{Using the Venn diagrams, if necessary, find the single sets equal to the following: $A\cup\phi$}
\sq{Using the Venn diagrams, if necessary, find the single sets equal to the following: $\phi\cup\phi$}
\sq{Consider the set $P=\{x\mid x=5m,m\in N\}$ and $Q=\{x\mid x=2m,m\in N\}$. Find $P\cap Q$}
\sq{If $g(x)=7x-2$ find: g(0)}
\sq{If $g(x)=7x-2$, then find: g(-1)}
\sq{If $g(x)=7x-2$ then, find: $g\!\left(-\dfrac{5}{3}\right)$}
\sq{If $s(x)=8x^{2}-3$ find: s(1)}
\sq{If $s(x)=8x^{2}-3$ find: s(-9)}
\sq{If $s(x)=8x^{2}-3$ find: $s\!\left(\dfrac{7}{2}\right)$}
\sq{Consider the function defined by $k(x)=7x-5$. if $k(x)=100$, find the value of x.}
\sq{Factorize by identifying common factors: 6x+12}
\sq{Factorize by identifying common factors: $15y^{2}+20y$}
\sq{Factorize by identifying common factors: $-12x^{2}-3x$}
\sq{Factorize by identifying common factors: $4a^{2}b+8ab^{2}$}
\sq{Factorize by identifying common factors: $xy-3x^{2}+2x$}
\sq{Factorize by identifying common factors: $3a^{2}b-9ab^{2}+15ab$}
\sq{Factorize and represent pictorially: 5x+15}
\sq{Factorize and represent pictorially: $x^{2}+4x+3$}
\sq{Factorize and represent pictorially: $x^{2}+6x+8$}
\sq{Factorize and represent pictorially: $x^{2}+4x+4$}
\sq{Factorize: $x^{2}+x-12$}
\sq{Factorize: $x^{2}+7x+10$}
\sq{Factorize: $x^{2}-6x+8$}
\sq{Factorize: $x^{2}-x-56$}
\sq{Factorize: $x^{2}-10x-24$}
\sq{Factorize: $y^{2}+4y-12$}
\sq{Factorize: $y^{2}+13y+36$}
\sq{Factorize: $x^{2}-x-2$}
\sq{Factorize: $2x^{2}+7x+3$}
\sq{Factorize: $4x^{2}+13x+3$}
\sq{Factorize: $3x^{2}+5x+2$}
\sq{Define common factor.}
\sq{Define quadratic expression.}
\sq{Find common factor of $x^{2}+2x$ concretely, pictorially and symbolically.}
\sq{Factorize $x^{2}-5x+4$ concretely, pictorially and symbolically.}
\sq{Factorize $x^{2}-3x-10$ concretely, pictorially and symbolically.}
\sq{Factorize: $x^{2}+9x+14$}
\sq{Factorize: $x^{2}-11x+24$}
\sq{Factorize: $p^{2}+11p+18$}
\sq{Factorize: $3x^{2}-4x-4$}
\sq{Factorize: $8x^{3}+12x^{2}+6x+1$}
\sq{Factorize: $27a^{3}+108a^{2}b+144ab^{2}+64b^{3}$}
\sq{Factorize: $x^{3}+48x^{2}y+108xy^{2}+216y^{3}$}
\sq{Factorize: $8x^{3}-125y^{3}+150xy^{2}-60x^{2}y$}
\sq{Factorize: $125a^{3}-1$}
\sq{Factorize: $64x^{3}+125$}
\sq{Factorize: $x^{6}-27$}
\sq{Factorize: $1000a^{3}+1$}
\sq{Factorize: $343x^{3}+216$}
\sq{Factorize: $27-512y^{3}$}
\sq{Factorize: $x^{3}-18x^{2}+108x-216$}
\sq{Factorize: $8x^{3}+27$}
\sq{Factorize: $x^{3}-27y^{3}$}
\sq{Find HCF by factorization method: $21x^{2}y,35xy^{2}$}
\sq{Find HCF by factorization method: $4x^{2}-9y^{2},2x^{2}-3xy$}
\sq{Find HCF by factorization method: $x^{3}-1,\ x^{2}+x+1$}
\sq{Find LCM of the following expressions by using prime factorization method: $2a^{2}b,\ 4ab^{2},6ab$}
\sq{Find LCM of the following expressions by using prime factorization method: $x^{2}+x,\ x^{3}+x^{2}$}
\sq{Find LCM of the following expressions by using prime factorization method: $a^{2}-4a+4,\ a^{2}-2a$}
\sq{Define HFC.}
\sq{Define LCM.}
\sq{Write relationship between HCF and LCM.}
\sq{Find the HCF of $6x^{2}y,\ 9xy^{2}$.}
\sq{Find the LCM of $4x^{2}y,\ 8x^{3}y^{2}$}
\sq{Find the square root of the following polynomials by factorization method: $x^{2}-8x+16$}
\sq{Find the square root of the following polynomials by factorization method: $9x^{2}+12x+4$}
\sq{Find the square root of the following polynomials by factorization method: $36a^{2}+84a+49$}
\sq{Find the square root of the following polynomials by factorization method: $64y^{2}-32y+4$}
\sq{Factorize the following: $4x^{3}+18x^{2}-12x$}
\sq{Factorize the following: $x^{3}+64y^{3}$}
\sq{Factorize the following: $-x^{2}-23x-60$}
\sq{Factorize the following: $x^{3}y^{3}-8$}
\sq{Factorize the following: $2x^{2}+7x+3$}
\sq{Factorize the following: $x^{4}+64$}
\sq{Find LCM and HCF by prime factorization method: $x^{2}+8x+16,\ x^{2}-16$}
\sq{Find LCM and HCF by prime factorization method: $x^{3}-9x,\ x^{2}-4x+3$}
\sq{Find square root by factorization and division method of the expression $16x^{4}+8x^{2}+1$.}
\sq{Factorize the following: $x^{4}+2x^{2}+9$}
\sq{Solve and represent the solution on a real line: 12x + 30 = -6}
\sq{Solve and represent the solution on a real line: $\dfrac{x}{3}+6=-12$}
\sq{Solve and represent the solution on a real line: $\dfrac{x}{2}-\dfrac{3x}{4}=\dfrac{1}{12}$}
\sq{Solve and represent the solution on a real line: 2=7 (2x + 4)+12x}
\sq{Solve following inequality and represent the solution on a real line: $x-6\leq-2$}
\sq{Solve following inequality and represent the solution on a real line: -9 $>$ -16 + x}
\sq{Solve following inequality and represent the solution on a real line: $3+2x\geq3$}
\sq{Solve following inequality and represent the solution on a real line: $6(x+10)\leq0$}
\sq{Shade the solution region for the following linear inequalities in xy-plane: $2x+y\leq6$}
\sq{Shade the solution region for the following linear inequalities in xy-plane: $3x+7y\geq21$}
\sq{Shade the solution region for the following linear inequalities in xy-plane: $3x-2y\geq6$}
\sq{Shade the solution region for the following linear inequalities in xy-plane: $5x-4y\leq20$}
\sq{Shade the solution region for the following linear inequalities in xy-plane: $2x+1\geq0$}
\sq{Shade the solution region for the following linear inequalities in xy-plane: $3y-4\leq0$}
\sq{Define linear equation.}
\sq{How to solve a linear equation in one variable?}
\sq{Solve the following equation and represent their solution on real line: 3x - 5 = 7}
\sq{Solve the following equation and represent their solution on real line: $\dfrac{x-2}{5}-\dfrac{x-4}{2}=2$}
\sq{Find solution of $\dfrac{2}{3}x-1<0$ and also represent it on a real line.}
\sq{Solve the inequality x + 2y $<$6.}
\sq{Solve the following linear inequalities in xy-plane. $2x\geq-3$}
\sq{Solve the following linear inequalities in xy-plane. $y\leq2$}
\sq{Define problem constraints.}
\sq{Define decision variables.}
\sq{Define feasible region.}
\sq{Define feasible solution.}
\sq{Define Non-negative constraints.}
\sq{Define corner point.}
\sq{Define vertex.}
\sq{Define objective function.}
\sq{Define optimal solution.}
\sq{Solve and represent their solutions on real line: $\dfrac{x+5}{3}=1-x$.}
\sq{Solve and represent their solutions on real line: $\dfrac{2x+1}{3}+\dfrac{1}{2}=1-\dfrac{x-1}{3}$.}
\sq{Solve and represent their solutions on real line: 3x + 7 $<$ 16}
\sq{Find in which quadrant the following angles lie. Write a co-terminal angle for the following: $65^{\circ}$}
\sq{Find in which quadrant the following angles lie. Write a co-terminal angle for the following: $135^{\circ}$}
\sq{Find in which quadrant the following angles lie. Write a co-terminal angle for the following: $-40^{\circ}$}
\sq{Find in which quadrant the following angles lie. Write a co-terminal angle for the following: $210^{\circ}$}
\sq{Find in which quadrant the following angles lie. Write a co-terminal angle for the following: $-150^{\circ}$}
\sq{Convert the following into degrees, minutes, and seconds: $123.456^{\circ}$}
\sq{Convert the following into degrees, minutes, and seconds: $58.7891^{\circ}$}
\sq{Convert the following into degrees, minutes, and seconds: $90.5678^{\circ}$}
\sq{Convert the following into decimal degrees: $65^{\circ}32'15''$}
\sq{Convert the following into decimal degrees: $42^{\circ}18'45''$}
\sq{Convert the following into decimal degrees: $78^{\circ}45'36''$}
\sq{Convert the following into radians: $36^{\circ}$}
\sq{Convert the following into radians: $22.5^{\circ}$}
\sq{Convert the following into radians: $67.5^{\circ}$}
\sq{Convert the following into degree: $\dfrac{\pi}{16}\,rad$}
\sq{Convert the following into degree: $\dfrac{11\pi}{5}\,rad$}
\sq{Convert the following into degree: $\dfrac{7\pi}{6}\,rad$}
\sq{Find the arc length and area of a sector if: r = 6 cm central angle $\theta=\dfrac{\pi}{3}$ radians.}
\sq{Find the arc length and area of a sector if: $r=\dfrac{4.8}{\pi}$ and central angle $\theta=\dfrac{5\pi}{6}$.}
\sq{Find the percentage of the area of sector subtending an angle $\dfrac{\pi}{8}$ radians.}
\sq{What is trigonometry?}
\sq{When an angle will be positive?}
\sq{When an angle will be negative?}
\sq{What are co-terminal angles?}
\sq{Define degree.}
\sq{Define radian.}
\sq{Define full circle.}
\sq{Define Half circle.}
\sq{Define right angle.}
\sq{Prove that: l=r$\theta$}
\sq{Prove that: $A=\dfrac{1}{2}r^{2}\theta$}
\sq{Convert $73.12^{\circ}$ to degrees, minutes and seconds.}
\sq{Convert $109.42^{\circ}$ to degrees, minutes and seconds.}
\sq{Convert $45^{\circ}45'45''$ to decimal degrees.}
\sq{Convert $94^{\circ}27'54''$ to decimal degrees.}
\sq{Convert radians to degree: $\dfrac{5\pi}{3}rad$}
\sq{Convert radians to degree: $\dfrac{7\pi}{6}rad$}
\sq{Convert radians to degree: $\dfrac{11\pi}{6}rad$}
\sq{Convert radians to degree: 1.2 rad}
\sq{Convert degree to radian: $15^{\circ}$}
\sq{Convert degree to radian: $75^{\circ}$}
\sq{Convert degree to radian: $315^{\circ}$}
\sq{Convert degree to radian: $15^{\circ}15'$}
\sq{Find the arc length of a sector with radius r = 10 cm and central angle $\theta=60^{\circ}$.}
\sq{Find the area of a sector with radius r = 8 cm and angle $\theta=45^{\circ}$.}
\sq{If arc length of a sector of radius 5 cm is 11 cm, find the angle subtended by the arc in radians and degrees.}
\sq{If $\theta$ lies in first quadrant, find the remaining trigonometric ratios of $\theta$: $\sin\theta=\dfrac{2}{3}$}
\sq{If $\theta$ lies in first quadrant, find the remaining trigonometric ratios of $\theta$: $\cos\theta=\dfrac{3}{4}$}
\sq{If $\theta$ lies in first quadrant, find the remaining trigonometric ratios of $\theta$: $\tan\theta=\dfrac{1}{2}$}
\sq{If $\theta$ lies in first quadrant, find the remaining trigonometric ratios of $\theta$: $\sec\theta=3$}
\sq{If $\theta$ lies in first quadrant, find the remaining trigonometric ratios of $\theta$: $\cot\theta=\sqrt{\dfrac{3}{2}}$}
\sq{Prove the following trigonometric identities: $(\sin\theta+\cos\theta)^{2}=1+2\sin\theta\cos\theta$}
\sq{Prove the following trigonometric identities: $\dfrac{\cos\theta}{\sin\theta}=\dfrac{1}{\tan\theta}$}
\sq{Prove the following trigonometric identities: $\cos^{2}\theta-\sin^{2}\theta=2\cos^{2}\theta-1$}
\sq{Prove the following trigonometric identities: $\cos^{2}\theta-\sin^{2}\theta=1-2\sin^{2}\theta$}
\sq{Show that: $(\sec^{2}\theta-1)\cos^{2}\theta=\sin^{2}\theta$}
\sq{Show that: $\tan\theta+\cos\theta=\sec\theta\,\mathrm{cosec}\,\theta$}
\sq{Find the value of the following trigonometric ratios without using the calculator: $\sin30^{\circ}$}
\sq{Find the value of the following trigonometric ratios without using the calculator: $\cos30^{\circ}$}
\sq{Find the value of the following trigonometric ratios without using the calculator: $\tan\dfrac{\pi}{6}$}
\sq{Find the value of the following trigonometric ratios without using the calculator: $\tan60^{\circ}$}
\sq{Find the value of the following trigonometric ratios without using the calculator: $\sec60^{\circ}$}
\sq{Find the value of the following trigonometric ratios without using the calculator: $\cos\dfrac{\pi}{3}$}
\sq{Find the value of the following trigonometric ratios without using the calculator: $\cot60^{\circ}$}
\sq{Find the value of the following trigonometric ratios without using the calculator: $\sin60^{\circ}$}
\sq{Find the value of the following trigonometric ratios without using the calculator: $\sec30^{\circ}$}
\sq{Find the value of the following trigonometric ratios without using the calculator: $\mathrm{cosec}\,30^{\circ}$}
\sq{Find the value of the following trigonometric ratios without using the calculator: $\sin45^{\circ}$}
\sq{Find the value of the following trigonometric ratios without using the calculator: $\cos\dfrac{\pi}{4}$}
\sq{Evaluate: $2\sin60^{\circ}\cos60^{\circ}$}
\sq{Evaluate: $2\cos\dfrac{\pi}{3}\sin\dfrac{\pi}{3}$}
\sq{Evaluate: $2\sin45^{\circ}+2\cos45^{\circ}$}
\sq{Evaluate: $\cos60^{\circ}\cos30^{\circ}-\sin60^{\circ}\sin30^{\circ}$}
\sq{Evaluate: $\sin60^{\circ}\cos30^{\circ}-\cos60^{\circ}\sin30^{\circ}$}
\sq{Evaluate: $\cos60^{\circ}\cos30^{\circ}+\sin60^{\circ}\sin30^{\circ}$}
\sq{Evaluate: $\tan\dfrac{\pi}{6}\cot\dfrac{\pi}{6}+1$}
\sq{Find the value of x,y and z from the following right angled triangle.
\tfig{\coordinate(A)at(0,0);\coordinate(B)at(4,0);\coordinate(C)at(4,2.3);
\draw(A)--(B)--(C)--cycle;\draw(3.75,0)--(3.75,.25)--(4,.25);\draw(.9,0)arc(0:30:.9);\node at(1.25,.2){$30^\circ$};
\node[left]at(A){$A$};\node[below right]at(B){$B$};\node[above]at(C){$C$};
\node[below]at(2,0){$y=4$ cm};\node[right]at(4,1.1){$x$};\node[above left]at(2,1.15){$z$};}}
\sq{Find the value of x,y and z from the following right angled triangle.
\tfig{\coordinate(A)at(0,0);\coordinate(B)at(3,0);\coordinate(C)at(3,3);
\draw(A)--(B)--(C)--cycle;\draw(2.75,0)--(2.75,.25)--(3,.25);\draw(.9,0)arc(0:45:.9);\node at(1.15,.3){$\frac{\pi}{4}$};
\node[below left]at(A){$A$};\node[below right]at(B){$B$};\node[above]at(C){$C$};
\node[below]at(1.5,0){$x$};\node[above left]at(1.5,1.5){$z$};\node[rotate=90,right]at(3,1.5){$y=\sqrt{3}$ cm};}}
\sq{Find the value of x,y and z from the following right angled triangle.
\tfig{\coordinate(A)at(0,0);\coordinate(B)at(4,0);\coordinate(C)at(4,4);
\draw(A)--(B)--(C)--cycle;\draw(3.75,0)--(3.75,.25)--(4,.25);\draw(.9,0)arc(0:45:.9);\node at(1.3,.3){$45^\circ$};
\draw([shift={(4,4)}]-90:.6)arc(-90:-135:.6);
\node[left]at(A){$A$};\node[below right]at(B){$B$};\node[above]at(C){$C$};
\node[below]at(2,0){$y=4$ cm};\node[right]at(4,2){$x$};\node[above left]at(2,2){$z$};}}
\sq{Find the value of x,y and z from the following right angled triangle.
\tfig{\coordinate(C)at(0,0);\coordinate(B)at(2,0);\coordinate(A)at(2,3.46);
\draw(C)--(B)--(A)--cycle;\draw(1.75,0)--(1.75,.25)--(2,.25);\draw(.7,0)arc(0:60:.7);\node at(1.05,.35){$60^\circ$};
\node[left]at(C){$C$};\node[below right]at(B){$B$};\node[above]at(A){$A$};
\node[below]at(1,0){$x$};\node[right]at(2,1.7){$y$};\node[rotate=60,above]at(1,1.73){$z=2$ cm};}}
\sq{Find the unknown side and angles of the following triangle.
\tfig{\coordinate(A)at(0,0);\coordinate(B)at(1.6,0);\coordinate(C)at(1.6,3);
\draw(A)--(B)--(C)--cycle;\draw(.5,0)arc(0:62:.5);\draw([shift={(C)}]-90:.5)arc(-90:-118:.5);
\node[left]at(A){$A$};\node[below right]at(B){$B$};\node[above]at(C){$C$};
\node[below]at(.8,0){$c=\sqrt{13}$ cm};\node[left,above]at(.8,1.5){$b$};\node[rotate=90,right]at(1.6,1.5){$a=\sqrt{3}$ cm};}}
\sq{Find the unknown side and angles of the following triangle.
\tfig{\coordinate(A)at(0,0);\coordinate(B)at(4,0);\coordinate(C)at(4,3.3);
\draw(A)--(B)--(C)--cycle;\draw(3.75,0)--(3.75,.25)--(4,.25);
\node[left]at(A){$A$};\node[below right]at(B){$B$};\node[above]at(C){$C$};
\node[below]at(2,0){$c=4$ cm};\node[above left]at(2,1.65){$b$};\node[rotate=90,right]at(4,1.65){$a=4$ cm};}}
\sq{Each side of a square field if 60 m long. Find the lengths of the diagonals of the field.}
\sq{Calculate the value of x in the following case:
\tfig{\coordinate(B)at(0,2.5);\coordinate(C)at(4.5,2.5);\coordinate(A)at(0,0);
\draw(B)--(C)--(A)--cycle;\draw(0,2.25)--(.25,2.25)--(.25,2.5);
\node[above left]at(B){$B$};\node[above right]at(C){$C$};\node[below]at(A){$A$};
\node[above]at(2.25,2.5){12 cm};\node[left]at(0,1.25){$x$ cm};\node[rotate=-29,above]at(2.25,1.25){20 cm};}}
\sq{Calculate the value of x in the following case:
\tfig{\coordinate(A)at(0,0);\coordinate(B)at(3,0);\coordinate(C)at(5,0);\coordinate(D)at(3,2.5);
\draw(A)--(C)--(D)--cycle;\draw(B)--(D);\draw(3,.25)--(3.25,.25)--(3.25,0);
\node[below left]at(A){$A$};\node[below]at(B){$B$};\node[below right]at(C){$C$};\node[above]at(D){$D$};
\node[above left,rotate=40]at(1.5,1.25){$x$ cm};\node[above right,rotate=-52]at(4,1.25){5 cm};\node[below]at(4,0){4 cm};
\node[below]at(3,-.45){$m\overline{AB}=m\overline{BC}$};}}
\sq{Solve triangle ABC, When $m\angle A=60^{\circ}$, $b=5$ cm, $m\angle B=90^{\circ}$}
\sq{Solve triangle ABC, when a = 2cm, $b=2\sqrt{2}$ cm and $m\angle B=90^{\circ}$.}
\sq{The angle of elevation of the top of a flag post from a point on the ground level 40 m away from the flag post is $60^{\circ}$. Find the height of the post.}
\sq{A tree is 72 m high. Find the angle of elevation of its top from a point 100 m away on the ground level.}
\sq{A ladder makes an angle of $60^{\circ}$ with the ground and reaches a height of 10m along the wall. Find the length of the ladder.}
\sq{Find the measure of an angle of elevation of the Sun, if a tower 300m high casts a shadow 450 m long.}
\sq{Convert the given angles from degrees to radians giving answer in term of $\pi$: $255^{\circ}$}
\sq{Convert the given angles from degrees to radians giving answer in term of $\pi$: $75^{\circ}45'$}
\sq{Convert the given angles from degrees to radians giving answer in term of $\pi$: $142.5^{\circ}$}
\sq{Convert the given angles from radians to degrees giving answer in degree and minutes: $\dfrac{17\pi}{24}$}
\sq{Convert the given angles from radians to degrees giving answer in degree and minutes: $\dfrac{11\pi}{16}$}
\sq{Prove the following identity: $\dfrac{\sin\theta}{1-\cos\theta}=\dfrac{1+\cos\theta}{\sin\theta}$}
\sq{Describe the location in the plane of the point $P(x,y)$, for which: $x>0$}
\sq{Describe the location in the plane of the point $P(x,y)$, for which: $x>0$ and $y>0$}
\sq{Describe the location in the plane of the point $P(x,y)$, for which: $x=0$}
\sq{Describe the location in the plane of the point $P(x,y)$, for which: $y=0$}
\sq{Describe the location in the plane of the point $P(x,y)$, for which: $x\geq0$ and $y<0$}
\sq{Describe the location in the plane of the point $P(x,y)$, for which: $y=0,\ x=0$}
\sq{Describe the location in the plane of the point $P(x,y)$, for which: $|x|=|y|$}
\sq{Describe the location in the plane of the point $P(x,y)$, for which: $x\geq3$}
\sq{Describe the location in the plane of the point $P(x,y)$, for which: $y>2$}
\sq{Describe the location in the plane of the point $P(x,y)$, for which: x and y have opposite signs.}
\sq{Find the distance between the points: $A(6,7),\ B(0,-2)$}
\sq{Find the distance between the points: $C(-5,-2),D(3,2)$}
\sq{Find the distance between the points: $L(0,3),M(-2,-4)$}
\sq{Find the distance between the points: $P(-8,-7),Q(0,0)$}
\sq{Find in the following: The distance between the two given points.}
\sq{Find in the following: Midpoint of the line segment joining the two points: $A(3,1)\ ,\ B(-2,-4)$}
\sq{Find in the following: Midpoint of the line segment joining the two points: $A(-8,3)\ ,\ (2,-1)$}
\sq{Find in the following: Midpoint of the line segment joining the two points: $A\!\left(-\sqrt{5},-\dfrac{1}{3}\right),B\left(-3\sqrt{5},\ 5\right)$}
\sq{Is the following point is at a distance of 15 units from the origin? $(\sqrt{176},\ 7)$}
\sq{Is the following point is at a distance of 15 units from the origin? $(10,-10)$}
\sq{Is the following point is at a distance of 15 units from the origin? $(1,15)$}
\sq{Write breifly about the history of geometry.}
\sq{Define origin.}
\sq{Define coordinate axes.}
\sq{Define quardrants.}
\sq{Define two dimensional rectangular coordinate system.}
\sq{Write distance formula.}
\sq{Write mid point formula.}
\sq{Find the distance between the points: $A(5,6),B(5,-2)$}
\sq{Show that the points $A(-1,2),B(7,5)$ and $C(2,-6)$ are vertices of a right triangle.}
\sq{The point $C(-5,3)$ is the centre of a circle and $P(7,-2)$ lies on the circle. What is the radius of the circle?}
\sq{Find the midpoint of the line segment joining the points $A(2,3)$ and $B(8,7)$.}
\sq{Find the slope and inclination of the line joining the points: $(-2,4);(5,11)$}
\sq{Find the slope and inclination of the line joining the points: $(3,-2);(2,7)$}
\sq{Find the slope and inclination of the line joining the points: $(4,6);(4,8)$}
\sq{By means of slope, show that the following points lie on the same line: $A(-1,-3);B(1,5);C(2,9)$}
\sq{By means of slope, show that the following points lie on the same line: $P(4,-5);Q(7,5);R(10,15)$}
\sq{By means of slope, show that the following points lie on the same line: $L(-4,6);M(3,8);N(10,10)$}
\sq{By means of slope, show that the following points lie on the same line: $X(a,2b);Y(c,a+b);Z(2c-a,2a)$}
\sq{Find K so that the joining $A(7,3);B(k,-6)$ and the line joining $C(-4,5);D(-6,4)$ are: Parallel}
\sq{Find K so that the joining $A(7,3);B(k,-6)$ and the line joining $C(-4,5);D(-6,4)$ are: perpendicular}
\sq{Find an equation of : the horizontal line through: $(7,-9)$}
\sq{Find an equation of : the vertical line through $(-5,3)$}
\sq{Find an equation of : through $A(-6,5)$ having slope 7.}
\sq{Find an equation of : through $(8,-3)$ having slope 0.}
\sq{Find an equation of : through $(-8,5)$ having slope undefined.}
\sq{Find an equation of : through $(-5,-3)$ and $(9,-1)$.}
\sq{Find an equation of : y-intercept : -7 and slope: -5}
\sq{Find an equation of : x-intercept: -3 and y-intercept: 4}
\sq{Find an equation of : x-intercept: -9 and slope: -4.}
\sq{Define slope.}
\sq{Define x-intercept.}
\sq{Define y-intercept.}
\sq{Show that the points $A(-3,6),B(3,2)$ and $C(6,0)$ are colinear.}
\sq{Show that the triangle with vertices $A(1,1),B(4,5)$ and $C(12,-1)$ is a right triangle.}
\sq{Find an equation of line through the points (-2,1) and (6,-4).}
\sq{Write down an equation of the line which cuts the x-axis at (2,0) and y-axis at (0,-4).}
\sq{Find an equation of the line through the point P(2,3) which form an isosceles triangle with the coordinates axes in the first quadrant.}
\sq{Transform the equation $5x-12y+39=0$ into: slope intercept form.}
\sq{Transform the equation $5x-12y+39=0$ into: Two-intercept form.}
\sq{Transform the equation $5x-12y+39=0$ into: normal form}
\sq{Transform the equation $5x-12y+39=0$ into: point-slope form.}
\sq{Transform the equation $5x-12y+39=0$ into: two-point form.}
\sq{Transform the equation $5x-12y+39=0$ into: symmetric form.}
\sq{If the houses of two friends are represented by coordinates $(2,6)$ and $(9,12)$ on a gird. Find the straight line distance between their houses if the grid units represent kilometres?}
\sq{Consider a straight trail (represented by coordinate plane) that start at point $(5,7)$ and ends at point $(15,3)$. what is the coordinate of the midpoint?}
\sq{The start and end points of a race track are given by coordinates $(3,9)$ and $(9,13)$. What is the midpoint of the track?}
\sq{The coordinates of two points on a road are $A(3,4)$ and $B(7,10)$. Find the midpoint of the road.}
\sq{What is aviation?}
\sq{What is navigation?}
\sq{What is latitude?}
\sq{What is longitude.}
\sq{On a map town A is at coordinates (2,3) and Town B is at (-4,-1), what is the distance between the two towns?}
\sq{Suppose two cities, City A and City B, are represented by the coordinates (3,4) and (7,1) on a map. Find the straight-line distance between the two cities.}
\sq{An Engineer is building a bridge between two point on a riverbank. Suppose the coordinates of the two points where the bridge will start and end are (2,5) and (8,9). Find the coordinates of the midpoint, which will represent the centre of the bridge.}
\sq{A pilot need to travel from city A (50,60) to city B (120, 150). Determine the heading angle the plane should take relative to the east direction.}
\sq{Abdul Hadi is traveling from point A(Latitude $10^{\circ}$ N, Longitude $50^{\circ}$ E) to point B(Latitude $20^{\circ}$ N, Longitude $60^{\circ}$ E) . Find the midpoint of his journey in terms of latitude and longitude.}
\sq{A landscaper is designing a straight pathway from P(2,3) to Q (8,9). What is the length of the pathway?}
\sq{Find the distance between two points A(2,3) and B(7,8) on a coordinate plane.}
\sq{Find the midpoint of the line segment joining the points (4,-2) and (-6,3).}
\sq{Calculate the gradient (slope) of the line passing through the point (1,2) and (4,6).}
\sq{If two lines are parallel and one line has a gradient of $\dfrac{2}{3}$, what is the gradient of the other line?}
\sq{An airplane needs to fly from  city A to coordinates (12, 5) to city B at coordinates (8, -4). Calculate the straight-line distance between these two cities.}
\sq{In a landscaping project, the path starts at (2,3) and ends at (10,7). Find the midpoint.}
\sq{For a line with a gradient of -3 and a y-intercept of 2, write the equation of the line in: Slope-intercept.}
\sq{For a line with a gradient of -3 and a y-intercept of 2, write the equation of the line in: Point-slope from using the point (1,2)}
\sq{For a line with a gradient of -3 and a y-intercept of 2, write the equation of the line in: Two point form using the points (1,2) and (4,-7).}
\sq{For a line with a gradient of -3 and a y-intercept of 2, write the equation of the line in: Intercepts form.}
\sq{For a line with a gradient of -3 and a y-intercept of 2, write the equation of the line in: Symmetric form.}
\sq{For a line with a gradient of -3 and a y-intercept of 2, write the equation of the line in: Normal form.}
\sq{Write the converse, inverse and contrapositive of the following conditional: $\sim p\rightarrow q$}
\sq{Write the converse, inverse and contrapositive of the following conditional: $q\rightarrow p$}
\sq{Write the converse, inverse and contrapositive of the following conditional: $\sim q\rightarrow\sim p$}
\sq{Write the converse, inverse and contrapositive of the following conditional: $\sim q\rightarrow\sim p$}
\sq{Write the truth table of the following: $\sim(p\vee q)\vee(\sim q)$}
\sq{Write the truth table of the following: $\sim(\sim q\vee\sim p)$.}
\sq{Write the truth table of the following: $(p\ \vee\ q)\leftrightarrow(p\ \wedge\ p)$}
\sq{Differentiate between a mathematical statement and its proof. Give two examples.}
\sq{What is the difference between an axiom and a theorem? Give examples of each.}
\sq{What is the importance of logical reasoning in mathematical proofs? Give an example to illustrate your point.}
\sq{Indicate whether it is an axiom, conjecture or theorem and explain your reasoning. There is exactly one straight line through any two points.}
\sq{Indicate whether it is an axiom, conjecture or theorem and explain your reasoning. Every even number greater than 2 can be written as the sum of two prime numbers.}
\sq{Indicate whether it is an axiom, conjecture or theorem and explain your reasoning. The sum of the angles in a triangle is $180^{o}$.}
\sq{Suppose x is an integer. Then x is odd if and only if $9x+4$ is odd.}
\sq{Suppose x is an integer. If x is odd, then $7x+5$ is even.}
\sq{If x and y are positive real numbers and $x^{2}<y^{2}$ then $x<y$.}
\sq{The sum of the interior angles of a triangle is $180^{o}$.}
\sq{If a,b and c are non-zero real numbers, prove that: $\dfrac{a}{b}\ \dfrac{c}{d}\Leftrightarrow ad=bc$}
\sq{If a,b and c are non-zero real numbers, prove that: $\dfrac{a}{b}\cdot\dfrac{c}{d}=\dfrac{ac}{bd}$}
\sq{If a,b and c are non-zero real numbers, prove that: $\dfrac{a}{b}+\dfrac{c}{b}=\dfrac{a+c}{b}$}
\sq{Define induction.}
\sq{Define deduction.}
\sq{Define statement.}
\sq{What is negation.}
\sq{What is conjuction.}
\sq{What is disjunction.}
\sq{Define equivalence.}
\sq{Define theorem.}
\sq{What is conjecture?}
\sq{Define axiom.}
\sq{What is Euclid axiom?}
\sq{What is Peano axiom?}
\sq{What is deductive prooof?}
\sq{Whether the following statement is true or false: Lahore is the capital of th Punjab and Quetta is the capital of Balochistan.}
\sq{Whether the following statement is true or false: $4<5\wedge8<10$}
\sq{Whether the following statement is true or false: $2+2=3\wedge6+6=10$}
\sq{10 is a positive integer or 0 is a rational number. Find truth value of this disjunction.}
\sq{Triangle can have two right angles or Lahore is the capital of Sindh. Find the truth Value of this disjunction.}
\sq{Prove that in any universal set, the empty set $\phi$ is a subset of any set A.}
\sq{Prove the following mathematical statement: If x is an odd integer, then $x^{2}$ is also an odd integer.}
\sq{Prove the following mathematical statement: The sum of two odd numbers is an even number.}
\sq{Find whether the solids are similar. All lengths are in cm.
\pfig{fig665}{5.5cm}}
\sq{In triangle ABC, the sides are given as $m\overline{AB}=6$ cm, $m\overline{BC}=9$ cm and $m\overline{CA}=12$ cm. In triangle DEF, the sides are given as $m\overline{DE}=10.5$ cm, $m\overline{EF}=15.75$ cm and $m\overline{FD}=21$ cm Prove that the triangles are similar. (image)}
\sq{Find the value of x in the following:
\pfig{fig667}{5cm}}
\sq{Find the value of x in the following:
\pfig{fig668}{5cm}}
\sq{Find the value of x in the following:
\pfig{fig669}{4.5cm}}
\sq{Find the value of x, y and z in the given figure:
\pfig{fig670}{4.5cm}}
\sq{If one pair of corresponding sides are parallel to each other, then the triangles so formed as shown in the figure are similar. in the figure, $\overline{AB}$ is parallel to $\overline{CD}$. (image)}
\sq{The perimeter of a regular octagon is 48 cm. Another octagon has sides that are 1.2 times the sides of the first octagon. What is the length of side of second octagon?}
\sq{Find the ratio of the areas of similar figures if the ratio of their corresponding lengths are: 1:3}
\sq{Find the ratio of the areas of similar figures if the ratio of their corresponding lengths are: 3:4}
\sq{Find the ratio of the areas of similar figures if the ratio of their corresponding lengths are: 2:7}
\sq{Find the ratio of the areas of similar figures if the ratio of their corresponding lengths are: 8:9}
\sq{The radii of two spheres are in the ratio 3:4. what is the ratio of their volumes?}
\sq{Two regular tetrahedrons have volumes in the ratio 8:27. What is the ratio of their sides?}
\sq{Two right cones have volumes in the ratio 64:125. What is the ratio of: their heights}
\sq{Two right cones have volumes in the ratio 64:125. What is the ratio of: their base areas?}
\sq{Find the missing value in the following similar solid:
\pfig{fig681}{6.5cm}}
\sq{Find the missing value in the following similar solid:
\pfig{fig682}{5.5cm}}
\sq{Find the missing value in the following similar solid:
\pfig{fig683}{6.5cm}}
\sq{Find the unknown volume in the following similar solid:
\pfig{fig684}{6.5cm}}
\sq{What is the sum of the interior angles of a decagon (10-sided polygon)?}
\sq{What is each exterior angle of a regular pentagon?}
\sq{Calculate the measure of each interior angle of a regular hexagon.}
\sq{If the sum of the interior angles of a polygon is $1260^{o}$, how many sides does the polygon have?}
\sq{In a parallelogram ABCD, $m\overline{AB}=10$ cm, $m\overline{AD}=6$cm and $m\angle BAD=45^{\circ}$, Calculate the length of diagonal $m\overline{AC}$.}
\sq{A rectangular wall is 10 m tall and 120 m wide. How many gallons of paint are needed to cover the wall, if one gallon covers 35 m$^{2}$?}
\sq{A rectangular floor is 12 m by 15 m. How many square tiles, each 1 m by 1 m, are needed to cover the floor?}
\sq{A rectangular wall has a length of 10 m and a width of 4 meters. If 1 litre of paint covers 7 m$^{2}$, how many liters of paint are needed to cover the wall?}
\sq{Find the area of the equilateral triangle ABC of side length s metres.}
\sq{Sketch the graph of the following linear function: $y=3x-5$}
\sq{Sketch the graph of the following linear function: $y=-2x+8$}
\sq{Sketch the graph of the following linear function: $y=0.5x-1$}
\sq{Plot the graph of the following function: $y=4^{x}$}
\sq{Plot the graph of the following function: $y=5^{-x}$}
\sq{Plot the graph of the following function: $y=\dfrac{1}{x-3}\ x\neq3$}
\sq{Plot the graph of the following function: $y=\dfrac{2}{x}+3,\ x\neq0$}
\sq{Plot the graph of the following function: $y=x^{\frac{1}{2}}$}
\sq{Plot the graph of the following function: $y=3x^{\frac{1}{3}}$}
\sq{Plot the graph of the following function: $y=2x^{-2}$}
\sq{What are graphs?}
\sq{Define quadratic function.}
\sq{Define cubic function.}
\sq{Define reciprocal function.}
\sq{Sketch the graph of: $y=2x-1$}
\sq{Plot the graphs of: $y=x^{2}$ and $y=-x^{2}$}
\sq{Sketch the graph of $y=2x^{2}-7x-9$ for $-3\leq x\leq6$.}
\sq{Plot the graph of the following cubic function for $-3\leq x\leq3$. $y=-x^{3}$}
\sq{Plot the graph of $y=2x^{3}-3x^{2}+x-5$ for $-2\leq x\leq3$.}
\sq{Sketch the graph of the following reciprocal function: $y=\dfrac{1}{x-0.5},x\neq0.5$}
\sq{Plot the graph of the exponential function: $y=2^{x},\ -6\leq x\leq6$.}
\sq{Plot the graph of the exponential function $y=e^{x}$.}
\sq{Plot the graph of: $y=x^{3}$ for $-3\leq x\leq3$.}
\sq{Plot the graph of: $y=x^{-1}$}
\sq{Plot the graph of: $y=2x^{\frac{1}{5}}$}
\sq{A company manufactures school bags. The function of producing x bags is C(x)= 1200 + 800x and the revenue from selling x bags is R(x)=50x. Find the break-even point.}
\sq{The demand and supply function for a product are given by the equation $P_{d}=400-5Q,P_{s}=3Q+24$: plot the graph of each function over the interval Q=0 to Q=300.}
\sq{A company manufactures school bags. The function of producing x bags is C(x)= 1200 + 800x and the revenue from selling x bags is R(x)=50x. Determine the profit or loss when 250 bags are sold.}
\sq{Shahid's salary S(x) in rupees is based on the following formula: S(x)= 45000 + 4500x, where x is the number of years he has been with the company. Sketch and interpret the graph of salary function for $0\leq x\leq5$.}
\sq{A company manufactures school bags. The function of producing x bags is C(x)= 1200 + 800x and the revenue from selling x bags is R(x)=50x. Plot the graph of both the functions and identify the break-even point.}
\sq{Plot the graph of the following function: $y=3^{-x}$ for x from -2 to 4.}
\sq{Plot the graph of the following function: $y=\dfrac{2}{x+7},x\neq-7$}
\sq{Plot the graph of the following for x from -5 to 5: $y=x^{2}-3$}
\sq{Plot the graph of the following for x from -5 to 5: $y=15-x^{2}$}
\sq{What is meant by locus?}
\sq{What is scalene triangle?}
\sq{What is isosceles triangle?}
\sq{What is equilateral triangle?}
\sq{What is acute angled triangle?}
\sq{What is obtuse angled triangle?}
\sq{What is right angled triangle?}
\sq{What is SSA?}
\sq{Define perpendicular bisector.}
\sq{Define point of concurrency.}
\sq{Define circumcenter.}
\sq{Define centeriod.}
\sq{What is angle of bisector of a triangle?}
\sq{Define Incenter.}
\sq{Define altitude.}
\sq{Define orthocenter.}
\sq{Define median.}
\sq{what is loci?}
\sq{Define circle.}
\sq{Define angle bisector.}
\sq{The following distribution represents the scores achieved by a group of chemistry students in the chemistry laboratory:
\dist What is the upper limit of the last class?}
\sq{The following distribution represents the scores achieved by a group of chemistry students in the chemistry laboratory:
\dist What is the lower limit of the class 39-43?}
\sq{The following distribution represents the scores achieved by a group of chemistry students in the chemistry laboratory:
\dist What is the midpoint of the class (34-38)?}
\sq{The following distribution represents the scores achieved by a group of chemistry students in the chemistry laboratory:
\dist What are the class frequencies of the classes 29-33 and 44-48?}
\sq{The following distribution represents the scores achieved by a group of chemistry students in the chemistry laboratory:
\dist What is the size of the class limits in the above frequency distribution?}
\sq{The following distribution represents the scores achieved by a group of chemistry students in the chemistry laboratory:
\dist In which class or group does minimum number of students fall?}
\sq{The following distribution represents the scores achieved by a group of chemistry students in the chemistry laboratory:
\dist What is the lower limit of the class having 15 as its class frequency?}
\sq{The following distribution represents the scores achieved by a group of chemistry students in the chemistry laboratory:
\dist What is the number of students having scores between 24 and 43?}
\sq{The following data show the numbers of heads in an experiment of 50 sets of tossing a coin 5 times. Make a discrete frequency distribution from the information:\\[2pt]
\hspace*{7em}3 3 4 0 5 4 3 3 1 2 4 5 0 3 2 4 4 0 0 0 5 5\\
\hspace*{7em}4 3 2 5 3 2 1 3 5 4 3 2 1 3 2 1 3 1 3 1 4 3}
\sq{Find the arithmetic mean in the following: 4,6,10,12,15,20,25,28,30}
\sq{Find the arithmetic mean in the following: 12,18,19,0,-19,-18,-12}
\sq{Find the arithmetic mean in the following: 6.5,11,12.3,9,8.1,16,18,20.5,25.}
\sq{Find the arithmetic mean in the following: 8,10,12,14,16,20,22.}
\sq{Following are the heights in (inches) of 12 students. Find the median height: 55,53,54,58,60,61,62,56,57,52,51,63.}
\sq{Following are the earning (in Rs) of ten workers: 88,70,72,125,115,95,81,90,95,90. Calculate: Arithmetic Mean.}
\sq{Following are the earning (in Rs) of ten workers: 88,70,72,125,115,95,81,90,95,90. Calculate: Median.}
\sq{Following are the earning (in Rs) of ten workers: 88,70,72,125,115,95,81,90,95,90. Calculate: Mode.}
\sq{Ten boys work on a petrol pump station. They get weekly wages as follows: Wages: (in Rs) 4250,4350,4400,4250,4350,4410,4500,4300,4500,4390. Find the arithmetic mean by short formula, median and mode of their wages.}
\sq{The arithmetic mean of 45 numbers is 80. Find their sum.}
\sq{Five numbers are 1,4,0,7,9. Find their mean, Median and mode.}
\sq{A set of data contains the values as 148,145,160,157,156,160. Show that Mode$>$Median$>$Mean.}
\sq{Calculate the arithmetic mean from the following information: $D=X-140,\ \sum D=500$ and n = 10}
\sq{Calculate the arithmetic mean from the following information: $U=\dfrac{x-130}{6},\ \sum U=-150$ and n = 15}
\sq{Calculate the arithmetic mean from the following information: if $D=x-25,\ \sum fD=300$ and $\sum f=20$}
\sq{Calculate the arithmetic mean from the following information: if $U=\dfrac{x-120}{5},\ \sum fU=60$ and $\sum f=100$}
\sq{For the following data, find the weighted mean:
\par\smallskip\noindent\hspace*{7em}\begin{tabular}{|l|c|c|}\hline \textbf{Item} & \textbf{Quantity} & \textbf{Cost of Item (in thousands)}\\ \hline Washing Machine & 5 & 35\\ Heater & 3 & 5\\ Stove & 2 & 13\\ Dispenser & 6 & 18\\ \hline\end{tabular}}
\sq{A company is planing its next year marketing budget across five years: yearly budget (in million) are: 5,7,8,6,7. Find the average budget for the next year.}
\sq{Ahmad obtained the following marks in a certain examination. Find the weighted mean if weights 5,4,2,3,2,4 respectively are allotted to the subjects.
\par\smallskip\noindent\hspace*{7em}\begin{tabular}{|c|c|c|c|c|c|}\hline \textbf{Urdu} & \textbf{English} & \textbf{Science} & \textbf{Math} & \textbf{Islamiyat} & \textbf{Computer}\\ \hline 78 & 65 & 80 & 90 & 85 & 72\\ \hline\end{tabular}}
\sq{The marks of a student in five examinations were 64,75,81,87,90. Find the arithmetic mean of the marks.}
\sq{Find the arithmetic mean using short formula for the runs made by a batsman: Runs: 40,45,50,52,50,60,56,70.}
\sq{The following are the scores made by batsman. Find the median of the data. 8,12,18,13,16,5,50.}
\sq{Following are the marks out of 100 obtained by 10 students in English. 23,15,35,48,41,5,8,9,11,51. Find the median of the data.}
\sq{The following data describes the marks of a student in different subjects and weights assigned to these subjects are also given:
\par\smallskip\noindent\hspace*{7em}\begin{tabular}{|l|c|c|c|}\hline Marks (x) & 74 & 78 & 74\\ \hline Weights (w) & 4 & 3 & 5\\ \hline\end{tabular}\par\smallskip\hspace*{7em}Find its weighted mean.}
\sq{Define the following: Frequency distribution.}
\sq{Define the following: Histogram (Unequal class limits).}
\sq{Define the following: Mean.}
\sq{Define the following: Median.}
\sq{Given below are marks obtained by 45 students in the monthly test of Biology:
\bioA With the reference of the above table find the following: Upper class boundary of the 5th class.}
\sq{Given below are marks obtained by 45 students in the monthly test of Biology:
\bioB With the reference of the above table find the following: Lower class boundaries of all the classes.}
\sq{Given below are marks obtained by 45 students in the monthly test of Biology:
\bioA With the reference of the above table find the following: Midpoint of all the classes.}
\sq{Given below are marks obtained by 45 students in the monthly test of Biology:
\bioB With the reference of the above table find the following: The class interval with the least frequency}
\sq{The marks of a student Saad in six test were 84,91,72,68,87,78. Find the arithmetic mean of his marks.}
\sq{Usman rolled a fair dice eight times. Each time their sum was recorded as 8,5,6,6,9,4,3,11. Find the median and mode of the sum.}
\sq{Arshad rolls a dice, with sides labelled L,M,N,O,P,U. What is the probability that the dice lands on consonant?}
\sq{Shazia throws a pair of fair dice. What will ne the probability of getting: Sum of dots is at least 4.}
\sq{Shazia throws a pair of fair dice. What will ne the probability of getting: product of both dots between 5 to 10.}
\sq{Shazia throws a pair of fair dice. What will ne the probability of getting: the difference between both the dots is equal to 4.}
\sq{Shazia throws a pair of fair dice. What will ne the probability of getting: Number at least 5 on the first dice and the number at least 4 on the scond dice.}
\sq{One alphabet is selected at random from the word " MATHEMATIC". Find the probability of getting: Vowel.}
\sq{One alphabet is selected at random from the word " MATHEMATIC". Find the probability of getting: Consonant.}
\sq{One alphabet is selected at random from the word " MATHEMATIC". Find the probability of getting: An E}
\sq{One alphabet is selected at random from the word " MATHEMATIC". Find the probability of getting: An A}
\sq{One alphabet is selected at random from the word " MATHEMATIC". Find the probability of getting: Not M}
\sq{One alphabet is selected at random from the word " MATHEMATIC". Find the probability of getting: Not T.}
\sq{What is probability?}
\sq{Write formula of probability.}
\sq{Define experiment.}
\sq{Define outcomes.}
\sq{Define favourable outcome.}
\sq{Define sample space.}
\sq{Define event.}
\sq{Define sample point.}
\sq{What is certain event?}
\sq{What is impossible event?}
\sq{What is likely event?}
\sq{What is unlikely event?}
\sq{What is equally likely events?}
\sq{Abdul Rehman rolls a fair dice, what is the probability of getting the number divisible by 3?}
\sq{If Zeeshan rolled two fair dice, find the probability of getting: Even numbers on both dice.}
\sq{If Zeeshan rolled two fair dice, find the probability of getting: Multiples of 3 on both dice.}
\sq{If Zeeshan rolled two fair dice, find the probability of getting: Even number on the first dice and the number 3 on the second dice.}
\sq{If Zeeshan rolled two fair dice, find the probability of getting: Al least the number 3 on the first dice and number 4 on the second dice.}
\sq{Zubair rolls a dice, what is the probability of not getting the number 6?}
\sq{If two fair dice are rolled. What is the probability of getting: not the sum of both dice is 8.}
\sq{If two fair dice are rolled. What is the probability of getting: not a double six.}
\sq{A card is drawn at random, from a pack of 52 playing cards. What is the probability of getting: a card of heart.}
\sq{A card is drawn at random, from a pack of 52 playing cards. What is the probability of getting: neither spade nor heart.}
\sq{In 500 trials of a thrown of two dice, what is expected frequency that the sum will be greater than 8?}
\sq{What is the expectation of a person who is toget-Rs. 120 if he obtains at least 2 heads in single toss of three coins?}
\sq{What is relative frequency?}
\sq{How relative frequency calculated?}
\sq{Define expected frequency.}
\sq{Find the average number of times getting 1 or 6, when a fair dice is rolled 300 times.}
\sq{If the probability of a defective bolt is 0.3. Find the numberyof non-defective bolts in a total to 800.}

\end{document}
